\documentclass[11pt,a4paper]{article}
\usepackage[utf8]{inputenc}
\usepackage[T1]{fontenc}
\usepackage{amsmath,amssymb,amsthm}
\usepackage{geometry}
\usepackage{fancyhdr}
\usepackage{microtype}
\usepackage{mathrsfs}

\geometry{margin=1in}

\pagestyle{fancy}
\fancyhf{}
\fancyhead[C]{\small Collected Mathematical Papers of Arthur Cayley, Vol.~IX}
\fancyfoot[C]{\thepage}
\renewcommand{\headrulewidth}{0.4pt}

\newcommand{\dd}{	ext{d}}

\begin{document}

\begin{center}
{\LARGE\bfseries COLLECTED MATHEMATICAL PAPERS\\[0.5em]
OF\\[0.5em]
ARTHUR CAYLEY, Sc.D., F.R.S.\\[1em]
VOL. IX.}
\end{center}

\vspace{2em}

\begin{center}
\rule{0.8\textwidth}{0.4pt}

\vspace{1em}
{\Large\scshape Pages 251--300}

\vspace{0.5em}
\rule{0.8\textwidth}{0.4pt}
\end{center}

\vspace{2em}

% ============================================================
% PAPER 596 -- pp. 251-252
% ============================================================
\section*{596.}
\begin{center}
{\large\bfseries ON JACOBI'S TRANSFORMATION OF A DOUBLE INTEGRAL.}\\[0.5em]
[From the \textit{Messenger of Mathematics}, vol.~IV. (1875), pp.~89--92.]
\end{center}

\textsc{Jacobi}, in the Memoir ``De Transformatione Integralis Duplicis
$\int\!\!\int A\,d\partial\,d\psi$ in formam $\int\!\!\int B\,dt\,d\theta$
transformanda,'' \textit{Crelle}, t.~XVI. (1837), pp.~21--25,
[\textit{Jacobi's Ges.\ Werke}, t.~III., pp.~91--97], considers the
transformation of
%
\begin{align*}
R^2 &= f^2(\sin\phi\cos\psi)^2(\sin\phi\sin\psi)^2
     + g^2(\sin\phi\cos\psi)^2(\cos\phi)^2
     + h^2(\cos\phi)^2(\sin\phi\cos\psi)^2 \\
    &\quad - a^2(\cos\phi)^2 - b^2(\sin\phi\cos\psi)^2 - c^2(\sin\phi\sin\psi)^2,\\[
\jot]
S^2 &= f^2(\sin\eta\cos\theta)^2(\sin\eta\sin\theta)^2
     + g^2(\sin\eta\sin\theta)^2(\cos\eta)^2
     + h^2(\cos\eta)^2(\sin\eta\cos\theta)^2 \\
    &\quad - a^2(\cos\eta)^2 - b^2(\sin\eta\cos\theta)^2 - c^2(\sin\eta\sin\theta)^2,
\end{align*}
%
and obtains the theorem:

If
\[
\frac{\sin\phi\,d\phi\,d\psi}{R}
  = \frac{\sin\eta\,d\eta\,d\theta}{S},
\]
then we have
\begin{align*}
f^2 &= \alpha\beta''\gamma' - \alpha\beta'\gamma'', \\
g^2 &= \alpha'\beta\gamma'' - \alpha'\beta''\gamma, \\
h^2 &= \alpha''\beta'\gamma - \alpha''\beta\gamma',
\end{align*}
$\&$c.

And it may be added that the integral equations are, so to speak, a complete
integral of the differential relation; viz. in virtue of the identity
$af+bg+ch=0$, the differential relation contains really only four constants;
the integral relations contain the six constants
$\alpha:\alpha':\alpha'':\alpha'''$ and $\beta:\beta':\beta'':\beta'''$, or we
have two constants introduced by the integration.

The best form of statement is, in the first theorem, to write $x$, $y$ for
$\cos\eta$, $\sin\eta$ ($x^2+y^2=1$), $\xi$, $\eta$ for $\cos\theta$,
$\sin\theta$ ($\xi^2+\eta^2=1$), and similarly in the second theorem to
introduce the variables $x$, $y$, $z$ connected by $x^2+y^2+z^2=1$, and
$\xi$, $\eta$, $\zeta$ connected by $\xi^2+\eta^2+\zeta^2=1$; then in the
first theorem $d\xi\eta$, $d\theta$ represent elements of circular arc, and
in the second theorem $\sin\phi\,d\phi\,d\psi$ and
$\sin\eta\,d\eta\,d\theta$ represent elements of spherical surface, and
the theorems are:

\medskip
\noindent\textsc{I.}\quad
If $(x,y)$ are coordinates of a point on the circle $x^2+y^2=1$, and
$(\xi,\eta)$ coordinates of a point on the circle $\xi^2+\eta^2=1$, and if
$ds$, $d\sigma$ are the corresponding circular elements, then
\[
\frac{ds}{d\sigma}
= \frac{\sqrt{(ax^2+by^2-c)}}{\sqrt{(a\xi^2+b\eta^2-c)}}
\]
has for its complete integral $ax\xi+by\eta-c=0$.

\medskip
\noindent\textsc{II.}\quad
If $(x,y,z)$ are coordinates of a point on the sphere $x^2+y^2+z^2=1$, and
$(\xi,\eta,\zeta)$ coordinates of a point on the sphere
$\xi^2+\eta^2+\zeta^2=1$; and if $ds$, $d\sigma$ are the corresponding
spherical elements, and
\begin{align*}
\beta\gamma'-\beta'\gamma &= f, & \alpha\delta'-\alpha'\delta &= a, \\
\gamma\alpha'-\gamma'\alpha &= g, & \beta\delta'-\beta'\delta &= b, \\
\alpha\beta'-\alpha'\beta &= h, & \gamma\delta'-\gamma'\delta &= c,
\end{align*}
(whence $af+bg+ch=0$); and for shortness
\begin{align*}
R &= \sqrt{(f^2y^2z^2+g^2z^2x^2+h^2x^2y^2-a^2x^2-b^2y^2-c^2z^2)},\\
S &= \sqrt{(f^2\eta^2\zeta^2+g^2\zeta^2\xi^2+h^2\xi^2\eta^2
       -a^2\xi^2-b^2\eta^2-c^2\zeta^2)},
\end{align*}
then the differential relation
\[
\frac{\sin\phi\,d\phi\,d\psi}{R}
  = \frac{\sin\eta\,d\eta\,d\theta}{S}
\]
has for its complete integral the system
\begin{align*}
\alpha x\xi + \beta y\eta + \gamma z\zeta + \delta &= 0,\\
\alpha' x\xi + \beta' y\eta + \gamma' z\zeta + \delta' &= 0,
\end{align*}
where by complete integral is meant a system of two equations containing two
arbitrary constants.

% ============================================================
% PAPER 597 -- pp. 253-256
% ============================================================
\newpage
\section*{597.}
\begin{center}
{\large\bfseries ON A DIFFERENTIAL EQUATION IN THE THEORY OF ELLIPTIC\\
FUNCTIONS.}\\[0.5em]
[From the \textit{Messenger of Mathematics}, vol.~IV. (1875), pp.~110--113.]
\end{center}

The differential equation
\[
(2-Q^2)\frac{d^2Q}{dk^2} - 3Q\left(\frac{dQ}{dk}\right)^2
  = 3(1-k^2)\frac{dQ}{dk},
\]
considered ante, p.~69, [594, this volume, p.~244], belongs to a class of
equations transformable into linear equations of the second order, and
consequently is such that, knowing a particular solution, we can obtain the
general solution.

In fact, assuming
\[
z = \frac{1}{\sqrt{Q}} \exp.\left(-\frac{1}{2}\int\frac{Q}{1-k^2}\,dk\right),
\]
the equation becomes
\[
(2-Q^2)\frac{d^2z}{dk^2}
  = 3(1-k^2)\left\{
      3(1-k^2)\frac{dz}{dk}
      + 6z
      - 3(1-k^2)z\frac{Q'}{Q}
    \right\},
\]
viz. omitting the terms in $\left(\dfrac{dz}{dk}\right)^2$ which destroy each
other, and dividing by $3(1-k^2)$, this is
\[
\frac{(2-Q^2)}{3(1-k^2)}\frac{d^2z}{dk^2}
  = 3(1-k^2)\frac{dz}{dk} + 6z - 3(1-k^2)z\frac{Q'}{Q},
\]
or finally
\[
\frac{d^2z}{dk^2}
  + \frac{2}{Q}\frac{dQ}{dk}\frac{dz}{dk}
  + \frac{z}{(1-k^2)}\frac{d^2Q}{dk^2}
  + \frac{3z}{(1-k^2)}\frac{dQ}{dk} = 0.
\]

But knowing a particular value of $Q$ we have
\[
z = \exp.\left(-\frac{1}{2}\int\frac{Q}{1-k^2}\,dk\right),
\]
a particular value of $z$, and thence in the ordinary manner the general
value of $z$, giving the general value of $Q$.

The solution given in my former paper may be exhibited in a more simple form
by introducing, instead of $k$, the variable $\alpha$ connected with it by the
equation
\[
k^2 = \frac{\alpha^2(2+\alpha)}{1+2\alpha}.
\]
We have in fact, \textit{Fundamenta Nova}, p.~25, [Jacobi's \textit{Ges.\ Werke},
t.~I., p.~76],
\[
u = \frac{2+\alpha}{1+2\alpha}, \qquad v = \frac{2+\alpha}{\alpha(1+2\alpha)},\]
viz. these expressions of $u$, $v$ in terms of the parameter $\alpha$, are
equivalent to, and replace, the modular equation
\[u^4 - v^4 + 2uv(1-u^2v^2) = 0.\]
We thence obtain
\[
\frac{(1+2\alpha)^2}{\alpha^2(2+\alpha)^2} = \frac{u^2}{v^2},
\]
that is,
\[
uv = \frac{\alpha\sqrt{(1+2\alpha)}}{\sqrt{(2+\alpha)}}, \qquad
\frac{u}{v} = \frac{(2+\alpha)}{\alpha(1+2\alpha)},
\]
\[
Q'-Q = \frac{1}{uv}, \qquad
Q'+Q = 2uv(1+u^2v^2),
\]
and the particular solution, $Q = \dfrac{u}{v} + 2uv$, becomes
\[Q = \frac{1}{\alpha}\sqrt{(2+\alpha)(1+2\alpha)}.\]

Introducing into the differential equation $\alpha$ in place of $k$, this is
found to be
\[
\{4\alpha^2(2+\alpha)^2 - Q^2\}
  \left\{
    \frac{d^2Q}{d\alpha^2}
    + \frac{3(1+2\alpha)}{\alpha(2+\alpha)}\frac{dQ}{d\alpha}
  \right\}
  - 3Q\left(\frac{dQ}{d\alpha}\right)^2 = 0.
\]
But from this form it at once appears that it is convenient in place of
$\alpha$ to introduce the new variable $\beta$, $=\alpha+\dfrac{1}{\alpha}$; the
equation thus becomes
\[
4Q^2\frac{d^2Q}{d\beta^2}
  + \frac{dQ}{d\beta}(3+6\beta Q^2-Q^4) - 12 = (\beta^2-1)(\beta^2-9)Q^2,
\]
satisfied by $Q=\sqrt{(5+2\beta)}$; or, what is the same thing, writing
$5+2\beta=\gamma^2$, that is, $\beta=-\dfrac{5}{2}+\dfrac{\gamma^2}{2}$, the
equation becomes
\[
4Q^2\frac{d^2Q}{d\gamma^2}
  + \frac{1}{\gamma}
    \frac{dQ}{d\gamma}(3+6\gamma^2Q^2-\gamma^4Q^4) - 12\gamma^2
  = (\gamma^2-1)(\gamma^2-9)\gamma^2Q^2,
\]
satisfied by $Q=\gamma$.

\newpage
\noindent Writing here
\[
Q = \frac{1}{2}(\gamma^2-1)(\gamma^2-9)z^2,
\]
we have for $z$ the equation
\[
\frac{d^2z}{d\gamma^2}
  + \frac{3\gamma^2-14}{\gamma(\gamma^2-1)(\gamma^2-9)}\frac{dz}{d\gamma}
  + \frac{24\gamma^2}{\{(\gamma^2-1)(\gamma^2-9)\}^2}z = 0,
\]
satisfied by
\[
z = \left(\frac{\gamma^2-9}{\gamma^2-1}\right)^{\frac{1}{2}}.
\]
[In fact, this value gives
\[
\frac{dz}{d\gamma} = (\gamma^2-9)^{-\frac{1}{2}}(\gamma^2-1)^{-\frac{3}{2}},
\]
\[
\frac{d^2z}{d\gamma^2}
  = (-12\gamma^3+57\gamma^2+36)
    (\gamma^2-9)^{-\frac{3}{2}}(\gamma^2-1)^{-\frac{5}{2}},
\]
which verify the equation as they should do.]

Representing for a moment the differential equation by
\[A\frac{d^2z}{d\gamma^2} + B\frac{dz}{d\gamma} + Cz = 0\]
and putting $z=z_1\int y\,d\gamma$, then assuming $z=z_1\int y\,d\gamma$, we find
\[
Az_1\frac{dy}{d\gamma}
  + y\left(2A\frac{dz_1}{d\gamma}+Bz_1\right) = 0,
\]
that is,
\[
\frac{1}{y}\frac{dy}{d\gamma}
  + \frac{2}{z_1}\frac{dz_1}{d\gamma} + \frac{B}{A} = 0,
\]
viz.
\[
\frac{1}{y}\frac{dy}{d\gamma}
  + \frac{2}{z_1}\frac{dz_1}{d\gamma}
  + \frac{3\gamma^2-14}{\gamma(\gamma^2-1)(\gamma^2-9)} = 0,
\]
or
\[
\frac{1}{y}\frac{dy}{d\gamma}
  + \frac{2}{z_1}\frac{dz_1}{d\gamma}
  - \frac{3\gamma}{\gamma^2}
  - \frac{1}{2}\frac{2\gamma}{\gamma^2-1}
  + \frac{3}{2}\frac{2\gamma}{\gamma^2-9} = 0,
\]
whence, integrating
\[
\log y + 2\log z_1 - 3\log\gamma - \tfrac{1}{2}\log(\gamma^2-1)
  + \tfrac{3}{2}\log(\gamma^2-9) = 0,
\]
that is,
\[
y = \frac{1}{z_1^2}
    \frac{\gamma^3\sqrt{(\gamma^2-1)}}{(\gamma^2-9)^{\frac{3}{2}}}
  = \frac{(\gamma+1)(\gamma+3)}{(\gamma-3)^2},
\]
\[
\int y\,d\gamma
  = \gamma + \log\frac{(\gamma+1)^2(\gamma+3)^2}{(\gamma-3)^2}.
\]

\newpage
Hence, the general value of $z$ is
\[
z = \left(\frac{\gamma^2-9}{\gamma^2-1}\right)^{\frac{1}{2}}
    \left\{
      \gamma + K + \log\frac{(\gamma+1)^2(\gamma+3)^2}{(\gamma-3)^2}
    \right\},
\]
the constants of integration being $K$ and $\gamma_0$, or, what is the same
the corresponding value of $Q$ being
\[
Q = \frac{1}{2}(\gamma^2-1)(\gamma^2-9)z^2,
\]
which contains the single arbitrary constant $K$; when this vanishes, we have
the foregoing particular solution $Q=\gamma$.

I recall that the expression of $\gamma$ is
\begin{align*}
\gamma &= \sqrt{(5+2\beta)},
       = \sqrt{5+2\left(\alpha+\frac{1}{\alpha}\right)},\\
     &= \frac{1}{\sqrt{\alpha}}\sqrt{\{(2+\alpha)(1+2\alpha)\}},
\end{align*}
where $\alpha$ is connected with $k$ by the relation
\[
\frac{\alpha^2(2+\alpha)}{1+2\alpha} = k^2.
\]

% ============================================================
% PAPER 598 -- pp. 257-258
% ============================================================
\newpage
\section*{598.}
\begin{center}
{\large\bfseries NOTE ON A PROCESS OF INTEGRATION.}\\[0.5em]
[From the \textit{Messenger of Mathematics}, vol.~IV. (1875), pp.~149, 150.]
\end{center}

I had occasion to consider the integral
\[
\int_0^R \frac{r^{q-1}\,dr}{(r^2+e^2)^{\frac{1}{2}q+\delta}},
\]
where $e$ is small in regard to $R$ and $q$ is negative. The integral is
finite when $e=0$, and it might be imagined that it could be expanded in
positive powers of $e$; and, assuming it to be thus expansible, that the
process would simply be to expand under the integral sign in ascending
powers of $e$, and integrate each term separately, so that the series would
be in integer powers of $e^2$.

Take two particular cases. First, let $\delta=2$, $q=-\tfrac{4}{3}$; the
integral is
\[
\int_0^R \frac{r^{-\frac{7}{3}}\,dr}{(r^2+e^2)^{\frac{1}{3}}},
\]
which, putting $r=e\tan\theta$, becomes
\[
\frac{1}{e^{\frac{8}{3}}} \int_0^{\tan^{-1}\frac{R}{e}}
  \frac{\cos^{\frac{5}{3}}\theta\,d\theta}{\sin^{\frac{7}{3}}\theta},
\]
and this, putting $\cos\theta=t^{\frac{3}{2}}$, becomes
\[
\frac{3}{2e^{\frac{8}{3}}} \int \frac{t\,dt}{(1-t^3)^{\frac{5}{3}}},
\]
viz. the integral is not thus obtainable: the series is right as far as it
goes, but the true expansion contains a term in $e^2$; and the failure of
the series to give the true expansion is indicated by the appearance of
infinite coefficients. In fact, the indefinite integral is
$\tfrac{3}{2}(r^2+e^2)^{\frac{1}{3}}$; taking this between the limits, it is
$\ldots$

Again, let $\delta=1$, $q=-2$; the integral is
\[
\int_0^R \frac{r^{-3}\,dr}{(r^2+e^2)^{\frac{1}{2}}},
\]
which gives similarly
\[
\int_0^R \frac{r^{-3}\,dr}{(r^2+e^2)^{\frac{1}{2}}}
  = \int \left(r^{-4} - \tfrac{1}{2}e^2r^{-6} + \ldots\right)dr,
\]
viz. the integral is not thus obtainable: the series is right as far as it
goes, but the true expansion contains a term as $e^4\log e$, and the failure
is indicated by the infinite coefficients. In fact, the indefinite integral
is
\[
\left(\tfrac{1}{2}r^{-2} + \tfrac{1}{2}e^2r^{-4}\right)\sqrt{(r^2+e^2)}
  + \tfrac{3}{2}e^4\log\{r+\sqrt{(r^2+e^2)}\},
\]
which between the limits is
\[
\left(\tfrac{1}{2}R^{-2} + \tfrac{1}{2}e^2R^{-4}\right)\sqrt{(R^2+e^2)}
  + \tfrac{3}{2}e^4\log\frac{R+\sqrt{(R^2+e^2)}}{e}
  = \tfrac{1}{2}R^{-1} + \tfrac{1}{2}e^2R^{-2} + \ldots
  - \tfrac{3}{2}e^4\log e.
\]

In the general case, the term causing the failure is $Ke^{-\frac{2}{q}}$ when
$q$ is fractional, and $Ke^{-\frac{2}{q}}\log e$ when $q$ is integral.

As a step towards determining the entire expansion, I notice that, writing
$\omega = \dfrac{e}{r+e}$ or $r = e\omega(1-\omega)^{-1}$, the value of the
integral is
\[
= \frac{1}{e^{\delta}}\int_0^{\frac{R}{R+e}}
  \omega^{q-1}(1-\omega)^{\frac{1}{2}q-\delta}\,d\omega,
\]
where $\ldots$

% ============================================================
% PAPER 599 -- pp. 259-262
% ============================================================
\newpage
\section*{599.}
\begin{center}
{\large\bfseries A SMITH'S PRIZE DISSERTATION.}\\[0.5em]
[From the \textit{Messenger of Mathematics}, vol.~IV. (1875), pp.~157--160.]
\end{center}

\begin{center}
\textit{Write a dissertation on Bernoulli's Numbers and their use in Analysis.}
\end{center}

The function
\[
\frac{t}{e^t-1} + \frac{1}{2}t
\]
is an even function of $t$, as appears by expressing it in the form
\[
\frac{1}{2}t\frac{e^{\frac{1}{2}t}+e^{-\frac{1}{2}t}}{e^{\frac{1}{2}t}-e^{-\frac{1}{2}t}},
\]
and its value for $t=0$ being obviously $=1$, we may write
\[
\frac{t}{e^t-1} = 1 - B_1t^2 + B_2t^4 - B_3t^6 + \&\text{c.};
\]
or, what is the same thing,
\[
\frac{t}{e^t-1} = 1 - \frac{B_1}{1.2}t^2 + \frac{B_2}{1.2.3.4}t^4
  - \frac{B_3}{1.2.3.4.5.6}t^6 + \&\text{c.},
\]
where the several coefficients $B_1$, $B_2$, $B_3$, $\&$c., are, as is at once
seen, rational fractions, and, as it may be shown, are all of them positive.
These numerical coefficients $B_1$, $B_2$, $B_3$, $\&$c., are called Bernoulli's
numbers.

There is no difficulty in calculating directly the first few terms; viz. we
have
\[
\left(1 + \tfrac{1}{2}t + \tfrac{1}{6}t^2 + \tfrac{1}{24}t^3
  + \tfrac{1}{120}t^4 + \ldots\right)
\left(1 - \tfrac{B_1}{1.2}t^2 + \tfrac{B_2}{1.2.3.4}t^4 - \ldots\right) = 1,
\]
that is
\[
1 + t\left(-\tfrac{1}{2}\right)
  + t^2\left(-\tfrac{1}{2}\cdot\tfrac{1}{2} + \tfrac{1}{6}\right)
  + t^3\left(\tfrac{1}{2}\cdot\tfrac{B_1}{2} - \tfrac{1}{2}\cdot\tfrac{1}{6}
    + \tfrac{1}{24}\right) + \ldots = 1,
\]
viz.
\[
B_1 = \frac{1}{6},\quad
1 - 3B_1 + B_2 = 0, \quad \&\text{c.},
\]
which is therefore
\[
B_1 = \frac{1}{6}, \quad B_2 = \frac{1}{30}, \quad B_3 = \frac{1}{42},
\]
and so a few more terms might have been found.

But a more convenient method is to express the numbers in terms of the
differences of $0^m$ by means of a general formula for the expansion of a
function of $e^t$, viz. this is
\[
\phi(e^t) = \phi(1+\Delta)e^{0.t},
\]
where
\[
e^{0.t} = 1 + 0.t + \frac{0^2.t^2}{1.2} + \frac{0^3.t^3}{1.2.3} + \ldots,
\]
and the $\phi(1+\Delta)$ is to be applied to the terms $0^0$, $0^1$, $0^2$,
$0^3$, $\&$c. We have thus
\[
\frac{t}{e^t-1}
= \frac{t}{e^t-1}\cdot\frac{\log e^t}{\log e^t}
= \frac{\log(1+\Delta)}{\Delta}e^{0.t},
\]
\[
\frac{t}{e^t-1}
= \frac{\log(1+\Delta)}{\Delta}
  \left(1 + 0.t + \frac{0^2.t^2}{1.2} + \ldots\right).
\]
We have, as may be at once verified,
\[
\frac{\log(1+\Delta)}{\Delta}
= 1 - \tfrac{1}{2}\Delta + \tfrac{1}{3}\Delta^2 - \ldots
  + \frac{2B_n}{1.2\ldots 2n}\Delta^{2n},
\]
and by what precedes, since the coefficient of every higher odd power of $t$
vanishes,
\[
\frac{\log(1+\Delta)}{\Delta}
= 1 - \tfrac{1}{2}\Delta + \tfrac{1}{3}\Delta^2 - \ldots;
\]
and then, by comparing the even powers of $t$, that is,
\[
(-)^{n-1}B_n = \frac{\log(1+\Delta)}{\Delta}0^{2n},
\]
\[
(-)^{n-1}B_n
= \left(1 - \tfrac{1}{2}\Delta + \tfrac{1}{3}\Delta^2 - \ldots
  + \frac{2B_n}{1.2\ldots 2n}\Delta^{2n}\right)0^{2n},
\]
the series for $\dfrac{\log(1+\Delta)}{\Delta}$ being stopped at this point since
$\Delta^{2n+1}0^{2n}=0$, $\&$c.

For instance, in the case $n=1$, we have
\[
B_1 = \left(1 - \tfrac{1}{2}\Delta + \tfrac{1}{6}\Delta^2\right)0^2
    = 0'^2 - \tfrac{1}{2}(1^2-0^2) + \tfrac{1}{6}(2^2-2.1^2+0^2)
    = -\tfrac{1}{2} + \tfrac{2}{3} = \tfrac{1}{6},
\]
as above.

\newpage
The formula shows, not only that $B_n$ is a rational fraction but that its
denominator is at most $=$ least common multiple of the numbers $2$, $3$, \ldots,
$2n+1$; the actual denominator of the fraction in its least terms is, however,
much less than this, there being as to its value a theorem known as Staudt's
theorem.

It does not obviously show that the Numbers are positive, or afford any
indication of the rate of increase of the successive terms of the series.
These last requirements are satisfied by an expression for $B_n$ as the sum
of an infinite numerical series, which expression is obtained by means of the
function $\cot\theta$, as follows:

We have
\[
\frac{te^{\frac{1}{2}t}}{e^t-1}
= \frac{\tfrac{1}{2}t(e^{\frac{1}{2}t}+e^{-\frac{1}{2}t})}
         {e^{\frac{1}{2}t}-e^{-\frac{1}{2}t}}
= 1 + \frac{B_1}{1.2}t^2 - \frac{B_2}{1.2.3.4}t^4 + \ldots,
\]
or, writing herein $t=2i\theta$ ($i=\sqrt{(-1)}$ as usual), this is
\[
\frac{i\theta\cos\theta}{\sin\theta}
= 1 - \frac{2^2B_1}{1.2}\theta^2
  + \frac{2^4B_2}{1.2.3.4}\theta^4 - \ldots.
\]
But we have
\[
\theta\cot\theta
= 1 - \frac{B_1}{1.2}\theta^2
  - \frac{B_2}{1.2.3.4}\theta^4 - \ldots;
\]
\[
\log\sin\theta
= \log\theta + \log\left(1-\frac{\theta^2}{\pi^2}\right)
  + \log\left(1-\frac{\theta^2}{2^2\pi^2}\right) + \ldots,
\]
and thence, by differentiation,
\[
\cot\theta
= \frac{1}{\theta}
  - \frac{2\theta}{\pi^2-\theta^2}
  - \frac{2\theta}{4\pi^2-\theta^2} - \ldots;
\]
hence
\[
\theta\cot\theta
= 1 - 2\theta^2\left(\frac{1}{\pi^2-\theta^2}
    + \frac{1}{4\pi^2-\theta^2} + \ldots\right),
\]
that is,
\[
B_n = \frac{2(1.2\ldots 2n)}{(2\pi)^{2n}}
      \left(\frac{1}{1^{2n}} + \frac{1}{2^{2n}} + \frac{1}{3^{2n}} + \ldots\right),
\]
showing first, that $B_n$ is positive, and next, that it rapidly increases
with $n$, viz. $n$ being large, we have
\[
B_n = \frac{2(1.2\ldots 2n)}{(2\pi)^{2n}},
\]
or, instead of $1.2\ldots 2n$ writing its approximate value
$\sqrt{(2\pi)}\cdot(2n)^{2n+\frac{1}{2}}e^{-2n}$, this is
\[
B_n = 4\sqrt{(n\pi)}\left(\frac{n}{e\pi}\right)^{2n}.
\]

\newpage
The result may of course be considered from the opposite point of view, as
giving a determination of the sum
$\dfrac{1}{1^{2n}} + \dfrac{1}{2^{2n}} + \dfrac{1}{3^{2n}} + \ldots$
in terms of Bernoulli's Numbers, assumed to be known, viz. we thus have
\[
\frac{1}{1^{2n}} + \frac{1}{2^{2n}} + \frac{1}{3^{2n}} + \ldots
= \frac{(2\pi)^{2n}}{2(1.2\ldots 2n)}B_n.
\]
For instance, $n=1$,
\[
\frac{1}{1^2} + \frac{1}{2^2} + \ldots
= \frac{2\pi^2}{2.1.2}B_1 = \frac{\pi^2}{6},
\]
and this is one and a good instance of the use of Bernoulli's Numbers in
Analysis.

Another and very important one is in the summation of a series, or say in
the determination of $S_{m,x} = u_0x + u_1x + \ldots + u_{x-1}$; viz.
starting from
\[
\frac{t}{e^t-1}
= 1 - \frac{B_1}{1.2}t^2 + \frac{B_2}{1.2.3.4}t^4 - \ldots
\]
and writing therein $t=d_x$, and therefore $\dfrac{1}{e^t-1}=\dfrac{1}{e^{d_x}-1}=S$,
$\dfrac{t}{e^t-1}=\dfrac{d_x}{e^{d_x}-1}$, and applying each side to a function
$u_x$ of $x$, we have
\[
S\,d_x\,u_x = u_x - \frac{B_1}{1.2}d_x^2\,u_x
  + \frac{B_2}{1.2.3.4}d_x^4\,u_x - \ldots,
\]
or taking the two sides each between the integer limits $a$, $x$,
\[
\sum_a^x d_x\,u_x = [u_x]_a^x - \frac{B_1}{1.2}[d_x^2\,u_x]_a^x
  + \frac{B_2}{1.2.3.4}[d_x^4\,u_x]_a^x - \ldots,
\]
where if $u_x$ is a rational and integral function the series on the
right-hand side is finite.

If for instance $u_x=x$, the equation is
\[
a+(a+1)+\ldots+(x-1) = \tfrac{1}{2}(x^2-a^2) - \tfrac{1}{2}(x-a),
\]
viz.
\[
\{1+2+\ldots+(x-1)\} - \{1+2+\ldots+(a-1)\}
= \tfrac{1}{2}(x^2-a^2) - \tfrac{1}{2}(x-a),
\]
which is right.

Applying the formula to the function $\log x$, we deduce theorems as to the
$\Gamma$-function; and it is also interesting to apply it to $\dfrac{1}{x}$.

The above is given as a specimen of what might be expected in an
examination: I remark as faults the omission to make it clear that $B_n$ is
a rational fraction; and the giving the series-formula as a formula for the
convenient calculation of $B_n$. The omission to give the first-mentioned
straightforward process of development strikes me as curious.

% ============================================================
% PAPER 600 -- p. 263
% ============================================================
\newpage
\section*{600.}
\begin{center}
{\large\bfseries THEOREM ON THE $n$th ROOTS OF UNITY.}\\[0.5em]
[From the \textit{Messenger of Mathematics}, vol.~IV. (1875), p.~171.]
\end{center}

If $n$ be an odd prime, and $\alpha$ an imaginary $n$th root of unity, then
\[
\frac{n}{1-\alpha} = 1 + \alpha + \alpha^2 + \ldots + \alpha^{n-1},
\]
for instance, verified at once by means of the equation
$1+\alpha+\alpha^2+\ldots+\alpha^{n-1}=0$:
\[
\frac{5}{1-\alpha^2} = 4\left(\frac{1}{1-\alpha^2} + \frac{1}{1-\alpha^3}\right),
\]
where the term in $()$ is
\[
\frac{\alpha(1-\alpha^2)+\alpha^2(1-\alpha^4)}{(1+\alpha^2)(1+\alpha^4)}
= \frac{\alpha+1+\alpha^3+\alpha^4}{1+\alpha^2+\alpha^3+\alpha^4} = 1;
\]
and so in other cases.

The formula shows, not only that $S_n$ is a rational function of $n$, but
that it is equal to $\dfrac{n}{2}$ for all values of $n$.

% ============================================================
% PAPER 601 -- pp. 264-265
% ============================================================
\newpage
\section*{601.}
\begin{center}
{\large\bfseries NOTE ON THE CASSINIAN.}\\[0.5em]
[From the \textit{Messenger of Mathematics}, vol.~IV. (1875), pp.~187, 188.]
\end{center}

A symmetrical bicircular quartic has in general on the axis two nodofoci
and four ordinary foci; viz. joining a nodofocus with either of the circular
points at infinity, the joining line is a tangent to the curve at the
circular point (and, this being a node of the curve, the tangent has there a
three-pointic intersection): and joining an ordinary focus with either of the
circular points at infinity, the joining line is at some other point a
tangent to the curve, viz. an ordinary tangent of two-pointic intersection.

In the case of the Cassinian, each circular point at infinity is a
fleflecnode (node with an inflexion on each branch); of the four ordinary
foci on the axis, one coincides with one nodofocus, another with the other
nodofocus, and there remain only two ordinary foci on the axis; the
so-called foci of the Cassinian are in fact the nodofoci, viz. each of these
points is by what precedes a nodofocus plus an ordinary focus, and the line
from either of these points to a circular point at infinity, qu\`{a} tangent at
a fleflecnode, has there a four-pointic intersection with the curve.

The analytical proof is very easy; writing the equation under the
homogeneous form
\[
\{(x-az)^2+y^2\} \{(x+az)^2+y^2\} - b^4z^4 = 0,
\]
then the so-called foci are the points $(x=az, y=0)$, $(x=-az, y=0)$; at
either of these, say the first of them, the line drawn to one of the
circular points at infinity is $x=az+iy$, and substituting this value in the
equation of the curve we obtain $z^4=0$, viz. the line is a tangent of
four-pointic intersection; this implies that there is an inflexion at the
point of contact on the branch touched by the line $x=az+iy$; and there is
similarly an inflexion at the point of contact on the branch touched by the
line $x=-az+iy$; viz. the circular point $x=iy$, $z=0$ is a fleflecnode;
and similarly the circular point $x=-iy$, $z=0$ is a fleflecnode.

To verify that there are no real points at infinity, observe that, writing
$x^2+y^2=r^2$, $x^2-y^2=r^2\cos 2\theta$, the equation becomes
\[
r^4 - 2c^2r^2\cos 2\theta + c^4 - b^4 = 0,
\]
and putting $r=\infty$, we must have $\cos 2\theta=\infty$, which is impossible.

% ============================================================
% PAPER 602 -- pp. 266-280
% ============================================================
\newpage
\section*{602.}
\begin{center}
{\large\bfseries ON THE POTENTIALS OF POLYGONS AND POLYHEDRA.}\\[0.5em]
[From the \textit{Proceedings of the London Mathematical Society}, vol.~VI. (1874--1875),
pp.~20--34. Read December 10, 1874.]
\end{center}

The problem of the attraction of polyhedra is treated of by Mehler,
\textit{Crelle}, t.~LXVI. pp.~375--381 (1866); but the results here obtained
are exhibited under forms, which are very different from his and which give
rise to further developments of the theory.

\subsection*{General Formulae for the Potentials of a Cone and a Shell.}

\smallskip\noindent\textbf{1.}\quad
The law of attraction is taken to be according to the inverse square of the
distance; and I commence with the general case of a cone standing upon any
portion of a surface $\Sigma$ as its base, and attracting a point at its
vertex, the cone being considered as a mass of density unity.

\smallskip\noindent\textbf{2.}\quad
Considering, in the first instance, an element of mass, the position of which
is determined by its distance $r$ from the vertex (or origin) and by two
angular coordinates defining the position of the radius vector $r$, then the
element is $=r^2\,dr\,d\omega$ (where $d\omega$ is the element of solid
angle, or surface of the unit-sphere), and the corresponding element of
potential is $=r\,dr\,d\omega$; whence
\[
V = \int\!\!\int r\,dr\,d\omega,
\]
which, integrating from $r=0$ to $r=$ its value at the surface, is
$=\tfrac{1}{2}\int r^2\,d\omega$, where $r$ now denotes the radius vector at a
point of the surface, being, therefore, a given function of the two angular
coordinates: and the remaining (double) integration

\newpage
\noindent is to be extended to all values of the angular coordinates belonging to a
position of $r$ within the conical surface which is the other boundary of the
attracting mass, or say over the spherical aperture of the cone.

\smallskip\noindent\textbf{3.}\quad
If the value of the radius vector at the surface is taken to be $mr$ ($m$ a
constant), then we have obviously
\[
V = \tfrac{1}{2}m^2 \int r^2\,d\omega;
\]
and hence also, writing $m+dm$ instead of $m$, we obtain, for the potential
of the portion of the shell lying between the similar and similarly situated
surfaces $\Sigma$ and $\Sigma'$, belonging to the parameters $m$ and $m+dm$
respectively, the value
\[
V = m\,dm \int r^2\,d\omega;
\]
this is $=2\dfrac{dm}{m}$ into the potential of the cone; and we thus see that
it is the same problem to determine the potential of the cone, and that of
the subtended portion of the indefinitely thin shell included between the two
surfaces.

\smallskip\noindent\textbf{4.}\quad
The same result may be arrived at as follows: the element of solid angle
$d\omega$ determines on the surface an element of surface $dS$, and if $dv$
be the corresponding normal thickness of the shell, then the element of mass
is $=dv\,dS$, and the element of potential is $=\dfrac{dv\,dS}{mr}$ ($mr$
being, as before, the radius vector at the surface). Take $\alpha$ the
complement of the inclination of the radius vector to the tangent plane---that
is, $\alpha$ the inclination of the radius vector to the normal, or, what is
the same thing, to the perpendicular from the origin on the tangent plane
(whence, also, if $mp$ be the length of this perpendicular, then
$\cos\alpha = \dfrac{mp}{mr}$, that is, $m\,p = mr\cos\alpha$).

\newpage
\noindent\textbf{5.}\quad
We have $dv\,dS = dm\,p\,d\omega$, $= mr\cos\alpha\,dm\,d\omega$; and the
element of potential is thus $=\cos\alpha\,dm\,d\omega$, whence integrating
from $m=0$ to $m=1$, we have
\[
V = \int\!\!\int r\cos\alpha\,dm\,d\omega,
\]
which, since $r$ varies as $m$, gives $V = \tfrac{1}{2}\int r^2\cos\alpha\,d\omega$;
and comparing with the original value, we have the relation
$\int r^2\,d\omega = \int r^2\cos\alpha\,d\omega$, that is,
$\int r^2\sin\alpha\,d\omega = 0$; which may be verified directly.

\subsection*{Cone on a plane base, and plane figure.}

\smallskip\noindent\textbf{6.}\quad
Suppose that the surface $\Sigma$ is a plane; the surface $\Sigma'$ is, of
course, a parallel plane. Taking here $mp$ for the perpendicular distance of
the plane $\Sigma$ from the origin, then, if $\delta$ be the infinitesimal
distance of the two planes from each other, we have $\delta = p\,dm$, that
is, $dm = \dfrac{\delta}{p}$; the potential of the cone is, as before,
$= \tfrac{1}{2}m^2 \int r^2\,d\omega$, and that of the plane figure, thickness
$\delta$, is $= \dfrac{\delta}{p} \int r^2\,d\omega$.

\smallskip\noindent\textbf{7.}\quad
Taking, for greater convenience, $m=1$, we have
\begin{align*}
\text{Potential of cone} &= \tfrac{1}{2}\int r^2\,d\omega,\\
\text{Do.\ of plane figure} &= \frac{\delta}{p}\int r^2\,d\omega,
\end{align*}
where $p$ is now the perpendicular distance of the plane from the vertex; or
if, as regards the plane figure, the infinitesimal thickness $\delta$ is
taken as unity, then
\[
\text{Potential of plane figure} = \frac{1}{p}\int r^2\,d\omega.
\]
In each case $r$ is the value of the radius vector corresponding to a point
of the plane figure which is the base of the cone, and the integration
extends over the spherical aperture of the cone.

\smallskip\noindent\textbf{8.}\quad
If the position of the radius vector is determined by the usual angular
coordinates, $\theta$ its inclination to the axis of $z$, and $\phi$ its
azimuth from the plane of $zx$---viz. if we have
\[
x = r\sin\theta\cos\phi, \quad
y = r\sin\theta\sin\phi, \quad
z = r\cos\theta;
\]
then, as is well-known, $d\omega = \sin\theta\,d\theta\,d\phi$, and the
integral $\int r^2\,d\omega$ is $= \int\!\!\int r^2\sin\theta\,d\theta\,d\phi$.
Taking the inclination of $p$ to the axes to be $\alpha$, $\beta$, $\gamma$
respectively, the equation of the plane which is the base of the cone is
\[
x\cos\alpha + y\cos\beta + z\cos\gamma = p;
\]
viz. we have
\[
r = \frac{p}{(\cos\alpha\cos\phi + \cos\beta\sin\phi)\sin\theta + \cos\gamma\cos\theta},
\]
and the integral $\tfrac{1}{2}\int r^2\,d\omega$ is therefore
\[
\tfrac{1}{2}p^2 \int\!\!\int
  \frac{\sin\theta\,d\theta\,d\phi}
       {\{ (\cos\alpha\cos\phi + \cos\beta\sin\phi)\sin\theta
          + \cos\gamma\cos\theta \}^2};
\]
and, in particular, if $p$ coincide with the axis of $z$, so that the
equation of the plane is $z=p$, then the integral is
\[
\tfrac{1}{2}p^2 \int\!\!\int \frac{\sin\theta\,d\theta\,d\phi}{\cos^2\theta}
  = \tfrac{1}{2}p^2 \int\!\!\int \sec^2\theta \sin\theta\,d\theta\,d\phi.
\]

\newpage
\noindent\textbf{9.}\quad
The integration in regard to $\theta$ can be at once performed; viz. in the
latter case we have
\[
\int \frac{\sin\theta\,d\theta}{\cos^2\theta} = \sec\theta;
\]
and in the former case, writing, as we may do,
\[
(\cos\alpha\cos\phi + \cos\beta\sin\phi)\sin\theta + \cos\gamma\cos\theta
= \psi\cos(\theta-N),
\]
then
\[
\int \frac{\sin\theta\,d\theta}{\psi^2\cos^2(\theta-N)}
= \frac{1}{\psi^2}
  \int \frac{\sin(\theta-N+N)\,d\theta}{\cos^2(\theta-N)}
= \frac{1}{\psi^2}
  \left\{
    \cos N \sec(\theta-N) + \sin N \log\tan\left(\tfrac{1}{4}\pi + \tfrac{1}{2}(\theta-N)\right)
  \right\}.
\]

\subsection*{Case of a Polyhedron or a Polygon.}

\smallskip\noindent\textbf{10.}\quad
Consider now the pyramid, vertex the origin $O$, standing on a polygonal
base. Letting fall from the vertex a perpendicular $OM$ on the base of the
pyramid, and drawing planes through $OM$ and the several vertices of the
polygon, we thus divide the pyramid into triangular pyramids; viz. $AB$ being
any side of the polygon, a component pyramid (or tetrahedron) will be $OMAB$,
vertex and base $MAB$, where $MO$ is a perpendicular at $M$ to the triangular
base $MAB$. And drawing through $MO$ a plane at right angles to $AB$, meeting
it in $D$ (viz. $MD$ is the perpendicular from $M$ on the base $AB$ of the
triangle), we divide the triangular pyramid into two pyramids $OMAD$, $OMBD$,
each having

\newpage
\noindent $OMD$ for a common triangular base, and the potentials of these in regard to
the point $O$ are equal, but of opposite signs; the potential of the
triangular pyramid $OMAB$ is equal to twice that of the pyramid $OMAD$; that
is, the potential of any triangular pyramid $OMAB$ in regard to the point $O$
depends upon that of the pyramid $OADM$; and (what is the same thing) the
potential of any plane polygon in regard to the point depends upon that of
the right-angled triangle $ADM$, situate as above in regard to the point $O$.

\smallskip\noindent\textbf{11.}\quad
I take $OM=h$, $MD=f$, $DA=g$; viz. supposing, as we may do, that the plane
of the triangle is parallel to that of $xy$, the point $M$ on the axis of
$z$, and the side $MD$ parallel to the axis of $x$, then $f$, $g$, $h$ will
be the coordinates of the point $A$.

\subsection*{Formulae for component triangular Pyramid, and Triangle.}

\smallskip\noindent\textbf{12.}\quad
Writing, as above, $x=r\sin\theta\cos\phi$, $y=r\sin\theta\sin\phi$,
$z=r\cos\theta$, and observing that $h$ is the perpendicular distance
originally called $p$, we have, for the potential of the pyramid,
\[
V = \tfrac{1}{2}h^2 \int d\phi\,(\sec\theta),
\]
where, $\phi$ being regarded as a given angle, the integral expression
$\sec\theta$ must be taken from $\theta=0$ to the value of $\theta$
corresponding to a point in the side $AD$. For any such point we have
$f=r\sin\theta\cos\phi$, $h=r\cos\theta$, that is, $\dfrac{f}{h}=\tan\theta\cos\phi$,
or the required value of $\theta$ is $=\tan^{-1}\dfrac{f}{h\cos\phi}$, and
consequently that of $\sec\theta$ is
\[
\frac{\sqrt{(f^2+h^2\cos^2\phi)}}{h\cos\phi}
= \frac{\sqrt{(f^2+h^2+h^2\tan^2\phi)}}{h\sec\phi}
= \frac{\sqrt{(f^2+h^2+h^2\tan^2\phi)}}{\sqrt{(1+\tan^2\phi)}};
\]
or, as this may also be written,
\[
\sqrt{(f^2+h^2)\cos^2\phi + h^2\sin^2\phi};
\]
hence
\[
V = \tfrac{1}{2}h \int d\phi\,
    \sqrt{(f^2+h^2)\cos^2\phi + h^2\sin^2\phi}.
\]

\newpage
\noindent\textbf{13.}\quad
The first term of the integral, writing therein for a moment $\tan\phi=u$,
is
\[
\int \frac{\sqrt{(f^2+h^2+f^2u^2)}\,du}{(1+u^2)\sqrt{(f^2+h^2+f^2u^2)}}.
\]
Hence, replacing $u$ by its value, we have
\[
V = \tfrac{1}{2}h
    \left\{
      g\tan^{-1}\frac{g}{\sqrt{(f^2+h^2+g^2)}}
      + f\log\frac{\sqrt{(f^2+h^2+g^2)}+g}{\sqrt{(f^2+h^2)}}
      - h\tan^{-1}\frac{g}{h}
    \right\},
\]
to be taken from $g=0$ to the value of $g$ corresponding to the point $A$;
viz. we have here $f=r\sin\theta\cos\phi$, $g=r\sin\theta\sin\phi$,
$h=r\cos\theta$, and thence $\tan\phi=\dfrac{g}{f}$, or $g=f\tan\phi$;
whence, writing for shortness $s=\sqrt{(f^2+g^2+h^2)}$ (viz. $s$ denotes the
distance $OA$), we have
\[
V = \tfrac{1}{2}h
    \left\{
      g\tan^{-1}\frac{g}{s} + f\log\frac{s+g}{\sqrt{(f^2+h^2)}}
      - h\tan^{-1}\frac{g}{h}
    \right\};
\]
or, observing that
\[
\frac{s+g}{\sqrt{(f^2+h^2)}} = \frac{\sqrt{(f^2+h^2)}}{s-g},
\]
this is
\[
V = \tfrac{1}{2}h
    \left\{
      g\tan^{-1}\frac{g}{s} + \tfrac{1}{2}f\log\frac{s+g}{s-g}
      - h\tan^{-1}\frac{g}{h}
    \right\};
\]
for the potential of the pyramid $OMDA$ in regard to the point; by omitting
the factor $\tfrac{1}{2}h$, we have
\[
V = g\tan^{-1}\frac{g}{s} + \tfrac{1}{2}f\log\frac{s+g}{s-g}
  - h\tan^{-1}\frac{g}{h}
\]
for the potential of the triangle $MDA$. The expression $\tan^{-1}$ denotes,
here and elsewhere, an arc included between the limits $-\tfrac{1}{2}\pi$,
$+\tfrac{1}{2}\pi$; it is therefore $+$ or $-$ according as the tangent is
$+$ or $-$.

\newpage
\subsection*{Formula for rectangular Pyramid, and Rectangle.}

\smallskip\noindent\textbf{14.}\quad
Completing the rectangle $MDAE$, the potential of the triangle $MAE$ is
obtained by interchanging the letters $f$ and $g$; viz. we have
\[
V = f\tan^{-1}\frac{f}{s} + \tfrac{1}{2}g\log\frac{s+f}{s-f}
  - h\tan^{-1}\frac{f}{h}
\]
for the potential of the triangle $MEA$.

The sum of the two gives the potential of the rectangle $MDAE$; viz. for
this rectangle, we have
\[
V = g\tan^{-1}\frac{g}{s} + f\tan^{-1}\frac{f}{s}
  + \tfrac{1}{2}f\log\frac{s+g}{s-g} + \tfrac{1}{2}g\log\frac{s+f}{s-f}
  - h\tan^{-1}\frac{g}{h} - h\tan^{-1}\frac{f}{h}.
\]
But we have
\[
\tan^{-1}\frac{f}{s} + \tan^{-1}\frac{g}{s} + \tan^{-1}\frac{h}{s}
= \tfrac{1}{2}\pi;
\]
for the function on the left hand is
\[
\tan^{-1}\frac{f}{s} + \tan^{-1}\frac{g}{s} + \tan^{-1}\frac{h}{s}
= \tan^{-1}
  \frac{\dfrac{f}{s}+\dfrac{g}{s}+\dfrac{h}{s}
        -\dfrac{fgh}{s^3}}
       {1-\dfrac{fg+gh+hf}{s^2}},
\]
viz. the denominator being $1-\dfrac{fg+gh+hf}{f^2+g^2+h^2}=0$, the tangent of
the arc is $\infty$, and the component arcs being each positive and less than
$\tfrac{1}{2}\pi$, the arc is $=\tfrac{1}{2}\pi$.

We thus obtain
\[
V = g\tan^{-1}\frac{g}{s} + f\tan^{-1}\frac{f}{s}
  + \tfrac{1}{2}f\log\frac{s+g}{s-g} + \tfrac{1}{2}g\log\frac{s+f}{s-f}
  - \tfrac{1}{2}\pi h + h\tan^{-1}\frac{h}{s},
\]
or, substituting for $\tan^{-1}\dfrac{h}{s}$ its value
$\tfrac{1}{2}\pi - \tan^{-1}\dfrac{f}{s} - \tan^{-1}\dfrac{g}{s}$, this is
\[
V = \tfrac{1}{2}f\log\frac{s+g}{s-g} + \tfrac{1}{2}g\log\frac{s+f}{s-f}
  + h\left(\tan^{-1}\frac{f}{s} + \tan^{-1}\frac{g}{s}\right)
\]
for the potential of the rectangle.

\newpage
\noindent\textbf{15.}\quad
The potential of a rectangular pyramid, vertex $O$, base the rectangle
($x_0$, $y_0$), ($x_1$, $y_0$), ($x_1$, $y_1$), ($x_0$, $y_1$), the sides being
parallel to the axes of $x$ and $y$, and the attracted point being the
origin, is at once obtained from the foregoing value; viz. writing
$f=x_1-x_0$, $g=y_1-y_0$, and denoting by $s$ the distance from the origin
of the corner ($x_1$, $y_1$, $h$), that is, $s=\sqrt{(x_1^2+y_1^2+h^2)}$, we
have four terms each of the form in question; or, reducing by means of the
formulae
\[
\tan^{-1}\frac{x_1}{s} + \tan^{-1}\frac{y_1}{s} + \tan^{-1}\frac{h}{s}
= \tfrac{1}{2}\pi, \quad \&\text{c.},
\]
we obtain the expression for the potential of the rectangular pyramid.

\smallskip\noindent\textbf{16.}\quad
Similarly, the potential of the rectangle, sides $2a$, $2b$, at the point
$(0,0,h)$ on the axis, is
\[
V = a\log\frac{s+b}{s-b} + b\log\frac{s+a}{s-a}
  + 2h\left(\tan^{-1}\frac{a}{s} + \tan^{-1}\frac{b}{s}\right),
\]
where $s=\sqrt{(a^2+b^2+h^2)}$; and so for other cases.

\subsection*{General Formula for the Potential of a Polyhedron.}

\smallskip\noindent\textbf{17.}\quad
The potential of any polyhedron at an external point may be found by
dividing it into tetrahedra, and summing their potentials. For a tetrahedron
with one vertex at the attracted point, the potential has been found above.

\smallskip\noindent\textbf{18.}\quad
For a cuboid with vertices $(\pm a, \pm b, \pm c)$, the potential at an
exterior point $(x,y,z)$ with $x>a$, $y>b$, $z>c$, is obtained by summing
the contributions from the eight corners, with appropriate signs. The result
may be written
\begin{align*}
V &= \sum_{\pm} \Bigg[
    \tfrac{1}{2}(x\mp a)(y\mp b)\log\left\{(z\mp c)+r_{\pm\pm\pm}\right\} \\
  &\quad + \tfrac{1}{2}(x\mp a)(z\mp c)\log\left\{(y\mp b)+r_{\pm\pm\pm}\right\} \\
  &\quad + \tfrac{1}{2}(y\mp b)(z\mp c)\log\left\{(x\mp a)+r_{\pm\pm\pm}\right\} \\
  &\quad - \tfrac{1}{2}(x\mp a)^2
      \tan^{-1}\frac{(y\mp b)(z\mp c)}{(x\mp a)r_{\pm\pm\pm}} \\
  &\quad - \tfrac{1}{2}(y\mp b)^2
      \tan^{-1}\frac{(x\mp a)(z\mp c)}{(y\mp b)r_{\pm\pm\pm}} \\
  &\quad - \tfrac{1}{2}(z\mp c)^2
      \tan^{-1}\frac{(x\mp a)(y\mp b)}{(z\mp c)r_{\pm\pm\pm}}
  \Bigg],
\end{align*}
where $r_{\pm\pm\pm}=\sqrt{(x\mp a)^2+(y\mp b)^2+(z\mp c)^2}$.

\newpage
\noindent\textbf{19.}\quad
For a point in the plane of a face, say $z=0$, the potential of the cuboid
reduces to a simpler form. For a rectangle in the plane, the potential at
$(x,y,0)$ is obtained by taking $c=0$ in the above, with proper attention to
the limiting values of the inverse tangent terms.

\smallskip\noindent\textbf{20.}\quad
For the rectangle ($x_0$, $y_0$), ($x_1$, $y_0$), ($x_1$, $y_1$), ($x_0$, $y_1$),
the potential at a point $(x,y)$ in its plane, with $x>x_1$, $y>y_1$, is
\begin{align*}
V &= (x_1-x)\log\frac{r_{11}+(y_1-y)}{r_{10}+(y_0-y)}
  + (x-x_0)\log\frac{r_{01}+(y_1-y)}{r_{00}+(y_0-y)} \\
  &\quad + (y_1-y)\log\frac{r_{11}+(x_1-x)}{r_{01}+(x-x_0)}
  + (y-y_0)\log\frac{r_{10}+(x_1-x)}{r_{00}+(x-x_0)},
\end{align*}
where $r_{ij}^2=(x-x_i)^2+(y-y_j)^2$.

\smallskip\noindent\textbf{21.}\quad
The potential of the cuboid at an external point may also be expressed as
\[
V = B(x_0,y_0,z_0) - B(x_1,y_0,z_0) - B(x_0,y_1,z_0)
  + B(x_1,y_1,z_0) - \ldots,
\]
where $B(a,b,c)$ denotes a function involving logarithms and inverse
tangents.

\smallskip\noindent\textbf{22.}\quad
In verification whereof, observe that the remaining terms are
$=xy\cdot\text{(expression involving only $z$ and $r$)}$.

\subsection*{Application to the Potentials of the Point, the Line, the Rectangle, and the Cuboid.}

\smallskip\noindent\textbf{27.}\quad
Take now $V$ to be the potential of a cuboid, and consider the case where
the point is exterior to the cuboid. The potential can be written as a sum
over the eight vertices, with appropriate signs. For the potential of the
point $(0,0,0)$, we have
\[
V = \frac{1}{r}.
\]

\smallskip\noindent\textbf{28.}\quad
For the line segment from $(0,0,0)$ to $(a,0,0)$, the potential at $(x,y,z)$
is
\[
V = \log\frac{a-x+r_1}{-x+r_0},
\]
where $r_0^2=x^2+y^2+z^2$, $r_1^2=(a-x)^2+y^2+z^2$.

\newpage
\noindent For the rectangle ($0\leq x\leq a$, $0\leq y\leq b$), the potential is
\[
V = x\log\frac{b+\sqrt{(x^2+b^2)}}{\sqrt{(x^2)}}
  + y\log\frac{a+\sqrt{(a^2+y^2)}}{\sqrt{(y^2)}}
  - a\log\frac{b+\sqrt{(a^2+b^2)}}{a}
  - b\log\frac{a+\sqrt{(a^2+b^2)}}{b},
\]
where, as before, the attracted point is $(x,y)$ with $x>a$, $y>b$.

\smallskip\noindent\textbf{29.}\quad
The potential of a polygon at a point in its plane may be found by dividing
the polygon into triangles and summing their potentials. For each triangle,
the potential is obtained by the formula above.

\smallskip\noindent\textbf{30.}\quad
For the cuboid, the complete formula is as given in Article~18 above. The
formula may be verified by differentiation; the negative gradient gives the
attraction in each direction.

% ============================================================
% PAPER 603 -- pp. 281-300
% ============================================================
\newpage
\section*{603.}
\begin{center}
{\large\bfseries ON THE POTENTIAL OF THE ELLIPSE AND THE CIRCLE.}\\[0.5em]
[From the \textit{Messenger of Mathematics}, vol.~IV. (1875), pp.~126--132.]
\end{center}

\subsection*{1.}

The potential of an ellipse at a point in its plane is a fundamental problem
in the theory of attractions. Let the ellipse be
\[
\frac{x^2}{a^2} + \frac{y^2}{b^2} = 1,
\]
and let the attracted point be $(0,c)$ on the minor axis. The potential is
\[
V = \int\!\!\!\int \frac{dx\,dy}{r},
\]
where $r^2=x^2+(y-c)^2$, the integral being over the area of the ellipse.

\subsection*{2.}

The density being taken unity, the potential is
\[
V = \int_{-b}^{b} \int_{-a\sqrt{1-y^2/b^2}}^{a\sqrt{1-y^2/b^2}}
    \frac{dx\,dy}{\sqrt{x^2+(y-c)^2}}.
\]

The integration with respect to $x$ gives
\[
\int \frac{dx}{\sqrt{x^2+(y-c)^2}}
  = \log\left\{x+\sqrt{x^2+(y-c)^2}\right\},
\]
and hence, writing $f^2=a^2-b^2$, $k^2=\dfrac{f^2}{a^2}$, and
$y=b\sin\phi$, we obtain an expression involving elliptic integrals.

\newpage
\subsection*{3.}

To accomplish this, writing
\[
A = a^2+b^2+c^2, \qquad
B = a^2+b^2-c^2, \qquad
C = 2bc,
\]
the potential becomes
\[
V = \int_{-\frac{\pi}{2}}^{\frac{\pi}{2}}
    \frac{2a\cos\phi\,d\phi}
         {\sqrt{a^2\cos^2\phi+(b\sin\phi-c)^2}}.
\]

\subsection*{4.}

This may be written
\[
V = 2a \int_{-\frac{\pi}{2}}^{\frac{\pi}{2}}
    \frac{\cos\phi\,d\phi}
         {\sqrt{A-B\cos^2\phi-C\sin\phi}}.
\]

\subsection*{5.}

The expression under the radical may be transformed by writing
\[
\cos\phi = \frac{1-t^2}{1+t^2}, \qquad
\sin\phi = \frac{2t}{1+t^2},
\]
and the integral becomes
\[
V = 4a \int_{-1}^{1}
    \frac{(1-t^2)\,dt}
         {(1+t^2)\sqrt{(1+t^2)^2A-(1-t^2)^2B-2Ct(1+t^2)}}.
\]

\subsection*{6.}

Simplifying the expression under the radical, we have
\begin{align*}
&(1+t^2)^2A - (1-t^2)^2B - 2Ct(1+t^2) \\
&\quad = (A-B)t^4 - 2Ct^3 + 2(A+B)t^2 - 2Ct + (A-B).
\end{align*}

\subsection*{7.}

This quartic in $t$ can be factorized, and the integral reduces to elliptic
form. The roots of the quartic determine the modulus and the amplitude of
the elliptic integrals.

\newpage
\subsection*{8.}

The potential of the ellipse at a point on the minor axis is expressible in
terms of the complete elliptic integrals $K$ and $E$, with a modulus depending
on the eccentricity and the position of the point.

\subsection*{9.}

For the circle, $b=a$, $k=0$, and the potential simplifies considerably.
For a point on the circumference, $c=a$, and
\[
V = 4a \int_0^{\frac{\pi}{2}} \frac{\cos\phi\,d\phi}{\sqrt{2-2\sin\phi}}
  = 4a \int_0^{\frac{\pi}{2}}
    \frac{\cos\phi\,d\phi}{\sqrt{2}\sqrt{1-\sin\phi}}.
\]

\subsection*{10.}

Writing $\sin\phi=1-2\sin^2\psi$, this becomes
\[
V = \frac{4a}{\sqrt{2}} \int_0^{\frac{\pi}{4}}
    \frac{2\cos 2\psi\,d\psi}{\sqrt{2}\sin\psi}
  = 2a \int_0^{\frac{\pi}{4}} \frac{\cos 2\psi\,d\psi}{\sin\psi},
\]
which is a logarithmic integral.

\subsection*{11.}

For the circle, the potential at a point $(0,c)$ on the axis, with $c<a$,
is given by a simpler formula. The potential at an internal point is
\[
V = 2\pi a - \text{(a term depending on the distance from the centre)}.
\]

\subsection*{12.}

For the ellipse, writing the equation in the form
\[
\frac{x^2}{a^2} + \frac{y^2}{a^2(1-e^2)} = 1,
\]
the potential at $(0,c)$ is
\[
V = 2\pi a\sqrt{1-e^2}\cdot F\left(e, \sin^{-1}\frac{c}{a}\right),
\]
where $F$ is an elliptic integral of the first kind.

\newpage
\subsection*{13.}

More precisely, introducing the notation of Legendre, the potential is
\[
V = \frac{4a'}{k}\left\{(1-\tfrac{1}{2}k^2)K - E\right\},
\]
where $a'=b=a\sqrt{1-e^2}$, and $k=e$ is the eccentricity.

\subsection*{14.}

For an external point, the formula involves incomplete elliptic integrals.
Writing $\lambda$ for the amplitude, we have
\[
V = \frac{4b}{k}\left\{
      (1-\tfrac{1}{2}k^2)F(k,\lambda) - E(k,\lambda)
    \right\},
\]
where $\sin\lambda = \dfrac{c}{a}$.

\subsection*{15.}

The general case, for an arbitrary point $(x,y)$ in the plane of the
ellipse, requires a more elaborate analysis. The potential depends on the
roots of a certain quartic equation, and involves elliptic integrals of both
the first and second kinds.

\subsection*{16.}

For the potential at $(x,y)$, we have
\[
V = \int\!\!\!\int \frac{dx'\,dy'}{\sqrt{(x-x')^2+(y-y')^2}},
\]
the integral being over the ellipse $\dfrac{x'^2}{a^2}+\dfrac{y'^2}{b^2}\leq 1$.

\newpage
\subsection*{17.}

By a linear transformation, this can be reduced to the potential of a circle
at a corresponding point. Writing $x'=a\xi$, $y'=b\eta$, the ellipse
becomes the circle $\xi^2+\eta^2\leq 1$, and the potential becomes
\[
V = ab\int\!\!\!\int \frac{d\xi\,d\eta}
                         {\sqrt{(x-a\xi)^2+(y-b\eta)^2}}.
\]

\subsection*{18.}

This integral can be evaluated by transforming to polar coordinates. Writing
$\xi=r\cos\theta$, $\eta=r\sin\theta$, we have
\[
V = ab\int_0^{2\pi}\int_0^1
    \frac{r\,dr\,d\theta}
         {\sqrt{(x-ar\cos\theta)^2+(y-br\sin\theta)^2}}.
\]

\subsection*{19.}

The integration with respect to $r$ gives
\begin{align*}
&\int_0^1 \frac{r\,dr}{\sqrt{A-2Br+Cr^2}} \\
&\quad = \frac{1}{C}\left\{
      \sqrt{A-2B+C} - \sqrt{A}
      + \frac{B}{\sqrt{C}}
        \log\frac{B-C+\sqrt{C}\sqrt{A-2B+C}}
                    {B+\sqrt{C}\sqrt{A}}
    \right\},
\end{align*}
where $A=x^2+y^2$, $B=ax\cos\theta+by\sin\theta$,
$C=a^2\cos^2\theta+b^2\sin^2\theta$.

\newpage
\subsection*{20.}

The remaining integration with respect to $\theta$ leads to elliptic
integrals. The complete expression for the potential is
\[
V = ab\int_0^{2\pi} \Phi(\theta)\,d\theta,
\]
where $\Phi$ involves square roots and logarithms.

\subsection*{21.}

For the case of the circle, $a=b$, the expression simplifies. The potential
at a point $(0,c)$ on the axis is
\[
V = 2\pi a\left\{\sqrt{1+\frac{c^2}{a^2}} - \frac{c}{a}\right\}
  = 2\pi\left(\sqrt{a^2+c^2}-c\right),
\]
for an external point.

\subsection*{22.}

For an internal point, $c<a$, the potential is
\[
V = 2\pi a\left(1-\frac{c}{2a}\right) = 2\pi a - \pi c.
\]

\subsection*{23.}

For the circle, the potential at any point in the plane depends only on the
distance from the centre. At the centre, $V=2\pi a$.

\subsection*{24.}

For the ellipse, the potential at the centre is
\[
V = 4aE(e),
\]
where $E(e)$ is the complete elliptic integral of the second kind.

\newpage
\subsection*{25.}

The potential at a focus is of special interest. For the ellipse with
$f^2=a^2-b^2$, the potential at the focus $(f,0)$ is
\[
V = \frac{4a}{e}E(e) - \frac{4a(1-e^2)}{e}K(e),
\]
where $e=f/a$ is the eccentricity.

\subsection*{26.}

For the circle, $e=0$, and this becomes indeterminate; taking the limit,
$V=2\pi a$, which agrees with the known result.

\subsection*{27.}

We may, by a different method, obtain the potential of the ellipse in the
form of a definite integral. Using the formula
\[
\frac{1}{r} = \int_0^{\infty} e^{-\lambda r}\,d\lambda,
\]
we have
\[
V = \int_0^{\infty} d\lambda \int\!\!\!\int e^{-\lambda r}\,dx\,dy.
\]

\subsection*{28.}

The inner integral, over the ellipse, can be evaluated. Writing
$r^2=x^2+(y-c)^2$, and transforming to elliptic coordinates, we obtain
\[
V = \int_0^{\infty} \frac{2\pi ab}{\sqrt{(1+\lambda^2a^2)(1+\lambda^2b^2)}}
    e^{-\lambda c}\,d\lambda.
\]

\newpage
\subsection*{29.}

For $c=0$, this gives
\[
V = 2\pi ab \int_0^{\infty}
    \frac{d\lambda}{\sqrt{(1+\lambda^2a^2)(1+\lambda^2b^2)}},
\]
which is a standard form for the complete elliptic integral. Writing
$\lambda=\dfrac{\tan\phi}{a}$, this becomes
\[
V = \frac{2\pi ab}{a} \int_0^{\frac{\pi}{2}}
    \frac{d\phi}{\sqrt{\cos^2\phi+\frac{b^2}{a^2}\sin^2\phi}}
  = 2\pi b \int_0^{\frac{\pi}{2}}
    \frac{d\phi}{\sqrt{1-e^2\sin^2\phi}}
  = 2\pi b K(e).
\]

\subsection*{30.}

For the potential at $(0,c)$ with $c>0$, we have from the general formula
\[
V = 2\pi ab \int_0^{\infty}
    \frac{e^{-\lambda c}\,d\lambda}
         {\sqrt{(1+\lambda^2a^2)(1+\lambda^2b^2)}}.
\]
Writing $a^2=b^2+f^2$, $k^2=f^2/a^2$, and $\lambda a=\tan\phi$, this
becomes
\[
V = \frac{2\pi ab}{a}
    \int_0^{\frac{\pi}{2}}
    \frac{e^{-\frac{c}{a}\tan\phi}\sec^2\phi\,d\phi}
         {\sqrt{(1+\tan^2\phi)(1+k^2\tan^2\phi)}}.
\]

\newpage
\subsection*{31.}

Simplifying,
\[
V = 2\pi b \int_0^{\frac{\pi}{2}}
    \frac{e^{-\frac{c}{a}\tan\phi}\,d\phi}
         {\sqrt{1-k^2\sin^2\phi}}.
\]
This integral is not expressible in terms of elementary functions, but can be
expanded in powers of $c/a$.

\subsection*{32.}

For the circle, $b=a$, $k=0$, and
\[
V = 2\pi a \int_0^{\frac{\pi}{2}}
    e^{-\frac{c}{a}\tan\phi}\,d\phi.
\]
This integral can be evaluated; writing $t=\tan\phi$, we have
\[
V = 2\pi a \int_0^{\infty} \frac{e^{-\frac{c}{a}t}\,dt}{1+t^2}.
\]

\subsection*{33.}

This is a known integral, and we have
\[
\int_0^{\infty} \frac{e^{-pt}\,dt}{1+t^2}
  = \mathrm{Ci}(p)\sin p - \mathrm{si}(p)\cos p,
\]
where $\mathrm{Ci}$ and $\mathrm{si}$ are the cosine and sine integrals.

For the present purpose, it is more convenient to use the direct evaluation.
For the circle, with external point $(0,c)$, $c>a$, the potential is
\[
V = 2\pi\left(\sqrt{a^2+c^2}-c\right),
\]
as may be verified by direct integration.

\newpage
\subsection*{34.}

If $c=0$, $a$ being $>0$, then $b=a$, $k=0$, and
\[
V = 2\pi a \int_0^{\frac{\pi}{2}} d\phi = \pi^2 a,
\]
which is the potential at the centre of a circle of radius $a$ with unit
density; since the above formula gives $V=2\pi a$ for $c=0$; and the correct
value is $V=2\pi a$, as stated.

\subsection*{35.}

Suppose, to fix the ideas, $f=1$, and consider the points $(0,c)$ and
$(a,0)$, which have equal potentials.

First, if $a>f$ (that is, $a>1$), then writing $k=\dfrac{f}{a}$, the
relation is
\[
\sqrt{a^2+c^2}-c = aE(k) - (1-k^2)aK(k),
\]
and we have
\[
2\pi a E(k) - 2\pi a(1-k^2)K(k) = 2\pi(\sqrt{a^2+c^2}-c),
\]
which determines $c$ in terms of $a$.

\vspace{2em}
\begin{center}
\rule{0.8\textwidth}{0.4pt}
\end{center}

\end{document}
